\documentclass[11pt,a4paper]{article}

\usepackage[margin=1in]{geometry}
\usepackage{times}
\usepackage{amsmath,amssymb}
\usepackage{booktabs}
\usepackage{graphicx}
\usepackage{hyperref}
\usepackage{url}
\usepackage{setspace}
\usepackage{enumitem}
\onehalfspacing

\hypersetup{colorlinks=true,linkcolor=black,citecolor=black,urlcolor=blue}

% Auto-generated by scripts/export_paper_numbers.py — do not edit by hand.
% Source: web/public/results/metrics.json
% Re-run: python scripts/export_paper_numbers.py

\newcommand{\MetricModelVersion}{1.0.0}
\newcommand{\MetricImgSize}{160}
\newcommand{\MetricSegThreshold}{0.4}
\newcommand{\MetricMinArea}{40}
\newcommand{\MetricNTrain}{534}
\newcommand{\MetricNVal}{134}
\newcommand{\MetricNTest}{112}
\newcommand{\MetricParamsM}{12.1}
\newcommand{\MetricOnnxMb}{11.7}
\newcommand{\MetricFpMb}{46.3}
\newcommand{\MetricFpTestDice}{0.686}
\newcommand{\MetricFpTestDiceCIlo}{0.619}
\newcommand{\MetricFpTestDiceCIhi}{0.750}
\newcommand{\MetricFpLesionDice}{0.751}
\newcommand{\MetricFpAuc}{0.939}
\newcommand{\MetricFpIoU}{0.618}
\newcommand{\MetricServedTestDice}{0.697}
\newcommand{\MetricServedTestDiceCIlo}{0.632}
\newcommand{\MetricServedTestDiceCIhi}{0.766}
\newcommand{\MetricServedLesionDice}{0.764}
\newcommand{\MetricServedLesionDiceCIlo}{0.705}
\newcommand{\MetricServedLesionDiceCIhi}{0.821}
\newcommand{\MetricServedIoU}{0.630}
\newcommand{\MetricServedAuc}{0.931}
\newcommand{\MetricServedSens}{0.867}
\newcommand{\MetricServedSpec}{0.841}
\newcommand{\MetricServedEce}{0.111}
\newcommand{\MetricNormalFP}{12/19}
\newcommand{\MetricNormalFPcount}{12}
\newcommand{\MetricNormalN}{19}
\newcommand{\MetricNormalFPrate}{0.632}
\newcommand{\MetricAnnotRate}{0.469}
\newcommand{\MetricCleanDice}{0.623}
\newcommand{\MetricFlaggedDice}{0.747}
\newcommand{\MetricDiceInflation}{0.063}
\newcommand{\MetricLesionCleanDice}{0.701}
\newcommand{\MetricLesionFlaggedDice}{0.799}
\newcommand{\MetricBusbraN}{1875}
\newcommand{\MetricBusbraDice}{0.714}
\newcommand{\MetricBusbraAuc}{0.638}
\newcommand{\MetricBreastN}{256}
\newcommand{\MetricBreastDice}{0.627}
\newcommand{\MetricBreastAuc}{0.721}
\newcommand{\MetricBusuclmStatus}{scored}
\newcommand{\MetricBusuclmN}{640}
\newcommand{\MetricBusuclmPatients}{38}
\newcommand{\MetricBusuclmDice}{0.386}
\newcommand{\MetricBusuclmLesionDice}{0.679}
\newcommand{\MetricBusuclmAuc}{0.780}
\newcommand{\MetricBusuclmNormalFP}{320/413}
\newcommand{\MetricTtaSpearman}{0.613}
\newcommand{\MetricDiceAtEightyCov}{0.778}
\newcommand{\MetricLeakValDelta}{-0.018}
\newcommand{\MetricVtwoCleanDice}{0.626}
\newcommand{\MetricVoneCleanDice}{0.623}
\newcommand{\MetricVtwoBusbraDice}{0.626}
\newcommand{\MetricVoneBusbraDice}{0.714}
\newcommand{\MetricVtwoBreastDice}{0.608}
\newcommand{\MetricVoneBreastDice}{0.627}
\newcommand{\MetricMeasDiamCorr}{0.697}
\newcommand{\MetricSizeDiscord}{0.171}


\title{Breast Ultrasound Lesion Segmentation:\\
an In-Browser Research Demo on BUSI with External Validation}
\author{Surabhi Fadnavis\\
\textit{High-school senior, Georgia}\\
\texttt{https://github.com/surabhif/Breast-ultrasound-lesion-segmentation}\\
DOI: \url{https://doi.org/10.5281/zenodo.23286597}}
\date{}

\begin{document}
\maketitle
\begin{center}
\small\textit{Research demo --- not for clinical use. Model version \MetricModelVersion.}
\end{center}

\begin{abstract}
We present a student research project and in-browser demonstration for breast ultrasound
lesion segmentation on the BUSI dataset. A ResNet-18 U-Net is trained with grouped
near-duplicate splits, exported to ONNX, and served as a dynamic INT8 model
(\MetricOnnxMb\,MB) that runs fully client-side. On the held-out BUSI test set the served model
reaches Dice~\MetricServedTestDice\ (95\% CI [\MetricServedTestDiceCIlo,~\MetricServedTestDiceCIhi]),
lesion-only Dice~\MetricServedLesionDice, and benign-vs-malignant ROC-AUC~\MetricServedAuc.
Frozen weights are evaluated without tuning on BUS-BRA (Dice~\MetricBusbraDice, AUC~\MetricBusbraAuc),
BrEaST (Dice~\MetricBreastDice, AUC~\MetricBreastAuc), and BUS-UCLM (all-image
Dice~\MetricBusuclmDice, lesion Dice~\MetricBusuclmLesionDice, AUC~\MetricBusuclmAuc;
normal FPs~\MetricBusuclmNormalFP). Caliper/annotation audits, a pre-registered
model-swap policy, and test-time uncertainty summaries are included. Normal images remain
difficult (non-empty masks on \MetricNormalFP\ normals). This is a portfolio research demo, not a
medical device.
\end{abstract}

\section{Introduction}

Breast ultrasound is widely used as an adjunct to mammography, especially in dense breasts.
Public research datasets such as BUSI~\cite{aldhabyani2020busi} enable students and
independent researchers to study lesion segmentation and triage-style classification, but
reported numbers are easy to over-interpret when splits leak near-duplicates, images contain
burned-in calipers, or models never leave the training distribution.

Student portfolio projects often stop at a training notebook. Clinically adjacent demos that
run in a browser are rarer, yet they force the same engineering constraints that matter for
honest reporting: a fixed preprocessing path, a quantized artifact small enough to download,
and metrics that match what the user actually runs. They also make failure modes visible---empty
masks on normals, brittle classification off-distribution, and optimistic BUSI scores on
annotated frames---in a way a table alone does not.

This project builds an end-to-end, reproducible pipeline: download and prepare BUSI, train a
compact ResNet-18 U-Net, export and quantize to ONNX, and ship a React web demo that runs
inference in the browser with onnxruntime-web. Metrics, failure galleries, and external
validation live beside the demo on GitHub Pages. The goals are educational clarity and honest
reporting---not clinical deployment.

\textbf{Contributions.}
(1)~An in-browser INT8 demo with training--serving parity (shared resize and morphology).
(2)~Grouped near-duplicate splits and a caliper/annotation audit on BUSI.
(3)~Pre-registered external validation of frozen model~\MetricModelVersion\ on BUS-BRA and BrEaST.
(4)~A written MODEL\_POLICY for when a retrained candidate may replace the served model.
(5)~Open citation/freeze scaffolding (CITATION.cff, Zenodo metadata, number-locked preprint
sources) with a minted software DOI (\url{https://doi.org/10.5281/zenodo.23286597}).

\section{Related work and positioning}

U-Net-style encoders remain a default for biomedical segmentation~\cite{ronneberger2015unet}.
Breast ultrasound datasets have expanded beyond BUSI to multi-center collections with patient
metadata and richer clinical fields, including BUS-BRA~\cite{gomezflores2024busbra},
BrEaST~\cite{pawlowska2024breast}, and BUS-UCLM~\cite{vallez2025busuclm}. Literature Dice
numbers vary widely with splits, resolutions, and whether normals are included; we therefore
treat published tables as context, not as a leaderboard to beat, and document protocol
differences in the companion site page \texttt{docs/LITERATURE\_COMPARISON.md}.

This work is deliberately scoped as a \emph{research demo}: single student author, CPU-friendly
training, browser inference, and explicit non-clinical disclaimers. We do not claim
state-of-the-art segmentation. We do claim that the numbers shown next to the demo are the
numbers of the served artifact, regenerated from committed JSON.

\section{Methods}

\subsection{Datasets and splits}

\textbf{BUSI}~\cite{aldhabyani2020busi} provides breast ultrasound images with lesion masks
and labels (benign, malignant, normal). The archive used here is obtained via a public
Hugging Face mirror of \emph{Dataset\_BUSI\_with\_GT} ($\approx$780 images). BUSI does not
provide patient identifiers. Splits therefore keep perceptual-hash near-duplicate clusters
together so near-identical frames cannot appear in both train and test. Held-out sizes:
train~\MetricNTrain, val~\MetricNVal, test~\MetricNTest. This is \emph{not} a true patient-level split.

\textbf{External sets} (CC~BY~4.0; scored after protocol registration, no fine-tuning):
BUS-BRA~\cite{gomezflores2024busbra} (\MetricBusbraN\ images, \texttt{n\_patients}${}=1064$ in the
committed JSON) and BrEaST~\cite{pawlowska2024breast} (\MetricBreastN\ images / patients).
BUS-UCLM~\cite{vallez2025busuclm} is the secondary set (\MetricBusuclmN\ images /
\MetricBusuclmPatients\ patients after excluding Doppler/combined frames; status
\MetricBusuclmStatus). UDIAT and BUSIS are not included (institutional licence / no-redistribution
constraints).

\subsection{Model and training}

Architecture: ResNet-18 encoder U-Net~\cite{ronneberger2015unet} with an auxiliary
benign-vs-malignant head ($\approx$\MetricParamsM\ M parameters). Images are resized to
\MetricImgSize$\times$\MetricImgSize\ with a shared half-pixel-center bilinear path used in training,
Python evaluation, and the browser. Training used paired flips, small rotations, scale-crop,
and brightness/contrast augmentation; loss combines pos-weighted BCE and soft Dice for
segmentation with BCE for classification on B/M only. Optimizer: AdamW with cosine schedule;
encoder freeze for the first few epochs; checkpoint selected on validation lesion Dice
(CPU-friendly recipe; seed~42; fold~0 primary run, with a five-fold CV summary retained for
diagnostics).

\subsection{Post-processing and served artifact}

Validation-tuned operating point (frozen for all reports): segmentation threshold
\MetricSegThreshold, minimum connected-component area \MetricMinArea\ pixels (4-connectivity),
classification threshold $0.5$. The training FP32 checkpoint is exported to ONNX and
dynamically quantized to UINT8 weights for serving
(\MetricOnnxMb\,MB vs \MetricFpMb\,MB FP32). The live demo loads
\texttt{web/public/models/v\MetricModelVersion/busi\_unet.onnx}. Served INT8 metrics are computed
with onnxruntime on CPU using the same postprocess as the browser.

\subsection{Browser demo}

The Vite/React TypeScript app runs onnxruntime-web in a Web Worker with self-hosted WASM,
optional threshold sliders, sample galleries, expert-vs-model overlays, BI-RADS context pages,
and Results pages backed by committed JSON. No ultrasound pixels leave the user's browser for
inference; the About page states that the site uses no analytics cookies for tracking. A
captioned walkthrough video and an auto-filled PDF report are shipped under
\texttt{web/public/}.

\subsection{External validation protocol and MODEL\_POLICY}

Protocol text lives in \texttt{docs/EXTERNAL\_VALIDATION\_PROTOCOL.md} and was committed before
Phase~1 scoring. External scoring uses identical preprocess/postprocess and the frozen INT8
weights. A candidate v2 may replace served v1 before the Dec~4 freeze only if it matches or
beats v1 on (a)~BUSI clean-subset Dice and (b)~held-out / external Dice under the fairness
rule below, and (c)~BUSI test AUC drops by at most $0.02$. Otherwise v1 remains served and v2
is published as an experiment only.

\paragraph{Phase~4 fairness.}
When v2 is trained on BUSI train plus a patient-grouped BUS-BRA train split, the BUS-BRA
held-out set is \emph{same-source held-out} for v2 (not external). We therefore score
v1.0.0 INT8 on the identical held-out case IDs and compare against that baseline, not against
v1's full-set BUS-BRA Dice. BrEaST remains the only truly external test. We report seed
mean${}\pm{}$SD and paired-bootstrap 95\% confidence intervals for each delta, decide
promotion on the seed mean, and---if the mean passes---promote the median-performing
passing seed by clean Dice (not the best).

\subsection{Additional analyses}

Heuristic caliper/text flags, clean-vs-flagged Dice comparisons, a short leakage ablation
(grouped vs random splits), Telea inpaint retrain experiments, offline flip-based TTA
uncertainty, and BrEaST measurement-agreement summaries (research sizing only) are reported
on the site and in committed JSON under \texttt{results/} and
\texttt{web/public/results/metrics.json}. Clinician outreach was considered early and
explicitly revoked; no reviewer acknowledgements appear on the site.

\section{Results}

All headline numbers below are macros generated from
\texttt{web/public/results/metrics.json} via \texttt{scripts/export\_paper\_numbers.py}
(\texttt{paper/numbers.tex}). Re-export and recompile after any metrics change;
\texttt{scripts/check\_paper\_numbers\_drift.py} fails on drift.

\subsection{BUSI held-out test (served INT8)}

\begin{table}[ht]
\centering
\caption{BUSI test metrics for served INT8 (v\MetricModelVersion) vs FP32 training checkpoint.}
\begin{tabular}{@{}lll@{}}
\toprule
Metric & FP32 & INT8 (served) \\
\midrule
Test Dice (all) & \MetricFpTestDice & \MetricServedTestDice\ [\MetricServedTestDiceCIlo,~\MetricServedTestDiceCIhi] \\
Lesion-only Dice & \MetricFpLesionDice & \MetricServedLesionDice\ [\MetricServedLesionDiceCIlo,~\MetricServedLesionDiceCIhi] \\
Test IoU & \MetricFpIoU & \MetricServedIoU \\
B vs M ROC-AUC & \MetricFpAuc & \MetricServedAuc \\
Sensitivity / Specificity & --- & \MetricServedSens\ / \MetricServedSpec \\
ECE & --- & \MetricServedEce \\
Normal non-empty masks & --- & \MetricNormalFP \\
\bottomrule
\end{tabular}
\end{table}

Quantization changes Dice only slightly (INT8 often within bootstrap noise of FP32) while
shrinking the download for browser use. Classification AUC on BUSI remains high
(\MetricServedAuc) but calibration is imperfect (ECE~\MetricServedEce). The normal-class false
positive rate (\MetricNormalFPrate) is a first-order usability failure for any ``run on a random
ultrasound'' demo: empty-ground-truth frames still receive non-empty masks on
\MetricNormalFP\ cases.

\subsection{Cleaning / caliper audit}

About \MetricAnnotRate\ of BUSI images were flagged for likely calipers or burned-in text by a
heuristic audit. Overall Dice on flagged test frames (\MetricFlaggedDice) exceeds clean frames
(\MetricCleanDice); inflation of all-minus-clean Dice is \MetricDiceInflation. Lesion-only Dice shows a
similar gap (flagged~\MetricLesionFlaggedDice\ vs clean~\MetricLesionCleanDice). Readers should treat
headline BUSI Dice as partly driven by easier annotated frames. Classification AUC on this
stronger model is similar on clean vs flagged subsets; earlier weaker baselines showed larger
gaps, so shortcut learning should be checked per recipe rather than assumed absent.

\subsection{External validation (frozen INT8)}

\begin{table}[ht]
\centering
\caption{External Dice and AUC for frozen served INT8 (no tuning).}
\begin{tabular}{@{}lrrr@{}}
\toprule
Dataset & $n$ & Dice & AUC \\
\midrule
BUSI (internal INT8) & \MetricNTest & \MetricServedTestDice & \MetricServedAuc \\
BUS-BRA & \MetricBusbraN & \MetricBusbraDice & \MetricBusbraAuc \\
BrEaST & \MetricBreastN & \MetricBreastDice & \MetricBreastAuc \\
BUS-UCLM & \MetricBusuclmN & \MetricBusuclmDice$^{\dagger}$ & \MetricBusuclmAuc \\
\bottomrule
\end{tabular}
\end{table}

Segmentation transfers moderately (BUS-BRA Dice~\MetricBusbraDice; BrEaST~\MetricBreastDice;
BUS-UCLM lesion Dice~\MetricBusuclmLesionDice). BUS-UCLM all-image Dice~\MetricBusuclmDice\ is
pulled down by normal false positives (\MetricBusuclmNormalFP). Classification AUC drops
substantially off-distribution (BUS-BRA~\MetricBusbraAuc; BrEaST~\MetricBreastAuc;
BUS-UCLM~\MetricBusuclmAuc), underscoring that BUSI triage scores are not portable without further work.
Subgroup Dice by BI-RADS and scanner on BUS-BRA is reported in the site Results JSON for
exploratory inspection; those rows were not used for threshold selection.

\subsection{MODEL\_POLICY outcome (inpaint retrain)}

A caliper-inpaint retrain (candidate v2) slightly improved clean-subset Dice
(\MetricVtwoCleanDice\ vs v1~\MetricVoneCleanDice) and did not harm AUC under the $0.02$ budget, but
\emph{failed} external Dice on BUS-BRA (\MetricVtwoBusbraDice\ vs \MetricVoneBusbraDice) and BrEaST
(\MetricVtwoBreastDice\ vs \MetricVoneBreastDice). The site therefore continues to serve
v\MetricModelVersion. Publishing the failed swap is intentional: it shows that cleaning training
images can help an internal clean subset while hurting transfer.

\subsection{Uncertainty, leakage, measurements}

Flip-based TTA disagreement correlates with error (Spearman $\rho\approx$\MetricTtaSpearman);
keeping the most certain 80\% of cases yields mean Dice~\MetricDiceAtEightyCov\ on the BUSI test
summary. This is agreement under flips---not a calibrated error probability. A short
grouped-vs-random leakage ablation on this schedule did not show large positive val inflation
(random$-$grouped $\approx$\MetricLeakValDelta); leakage risk remains conceptually important because
BUSI lacks patient IDs. On BrEaST, model vs expert longest-diameter correlation (mm,
exploratory) is \MetricMeasDiamCorr; disagreement rate on a $20$\,mm threshold proxy is
\MetricSizeDiscord. Imaging size is not pathological T-stage; these figures are research
diagnostics only.

\section{Limitations}

\begin{itemize}[leftmargin=*,itemsep=0.35em]
\item \textbf{Not for clinical use.} No clinician review; not a medical device; no diagnostic
  claims. Site pages that mention BI-RADS or surgical sizing are educational framing only.
\item BUSI lacks patient IDs; near-duplicate grouping is a proxy, not a guarantee against
  leakage.
\item High false-positive rate on normals (\MetricNormalFP).
\item External classification AUC is much lower than internal BUSI AUC.
\item Resolution is limited to \MetricImgSize$^2$ for CPU/browser practicality.
\item BUS-UCLM all-image Dice is low (\MetricBusuclmDice) largely from normal false positives
  (\MetricBusuclmNormalFP); lesion Dice \MetricBusuclmLesionDice\ is the fairer transfer figure.
\item Caliper detector is heuristic, not OCR; literature comparisons depend on heterogeneous
  protocols across papers.
\item Single primary training seed for the served model; CV folds are diagnostic, not an
  ensemble deployed in the browser.
\end{itemize}

\section{Ethics, licensing, and credit}

Code in this repository is released under the MIT license. BUSI redistribution terms remain
those of the original dataset providers; only a minimal cited sample ships in the web tree.
External demo samples use CC~BY sources with attribution. Personal information on the site is
limited to the author's name and ``high-school senior, Georgia''.

\section{AI-use disclosure}

Project by Surabhi Fadnavis. The code, analysis, text, and video were produced with AI tools (Cursor).

\section{Reproducibility and availability}

Code, metrics JSON, ONNX weights, and the GitHub Pages demo are in
\url{https://github.com/surabhif/Breast-ultrasound-lesion-segmentation}.
Citation metadata: \texttt{CITATION.cff}; Zenodo deposit metadata: \texttt{.zenodo.json}.
Software DOI: \url{https://doi.org/10.5281/zenodo.23286597}
(cite as~\cite{fadnavis2026busi}; agents must not create tags or releases). Paper numbers:
\begin{verbatim}
python scripts/export_paper_numbers.py
cd paper && pdflatex main && pdflatex main
python scripts/check_paper_numbers_drift.py
\end{verbatim}
Venue options (JEI vs arXiv vs site-only) are compared in \texttt{docs/SUBMISSION\_GUIDE.md};
nothing in this preprint package constitutes a submission.

\section*{Acknowledgments}

No clinician reviewers. Training used public datasets cited above. Hosting is GitHub Pages.

\bibliographystyle{plain}
\begin{thebibliography}{9}

\bibitem{fadnavis2026busi}
S.~Fadnavis.
\newblock Breast Ultrasound Lesion Segmentation: an in-browser research demo on BUSI with external validation (Version 1.0.0) [Computer software].
\newblock Zenodo, 2026.
\newblock \url{https://doi.org/10.5281/zenodo.23286597}.

\bibitem{aldhabyani2020busi}
W.~Al-Dhabyani, M.~Gomaa, H.~Khaled, and A.~Fahmy.
\newblock Dataset of breast ultrasound images.
\newblock \emph{Data in Brief}, 28:104863, 2020.
\newblock \url{https://doi.org/10.1016/j.dib.2019.104863}.

\bibitem{gomezflores2024busbra}
W.~G\'omez-Flores et al.
\newblock BUS-BRA: A breast ultrasound dataset for assessing computer-aided diagnosis systems.
\newblock \emph{Medical Physics}, 51:3110--3123, 2024.
\newblock \url{https://doi.org/10.1002/mp.16812}.

\bibitem{pawlowska2024breast}
A.~Paw{\l}owska et al.
\newblock A curated benchmark dataset for ultrasound-based breast lesion analysis (BrEaST).
\newblock \emph{Scientific Data}, 11:148, 2024.
\newblock \url{https://doi.org/10.1038/s41597-024-02984-z}.

\bibitem{vallez2025busuclm}
N.~Vallez et al.
\newblock BUS-UCLM: Breast ultrasound dataset from University of Castilla-La Mancha.
\newblock \emph{Scientific Data}, 12:242, 2025.
\newblock \url{https://doi.org/10.1038/s41597-025-04564-2}.

\bibitem{ronneberger2015unet}
O.~Ronneberger, P.~Fischer, and T.~Brox.
\newblock U-Net: Convolutional networks for biomedical image segmentation.
\newblock In \emph{MICCAI}, pages 234--241, 2015.

\end{thebibliography}

\end{document}
